\chapter{Stochastic satisfiability and inference under random choices}
\label{sec:p15-stochastic-satisfiability-and-inference}

\section{A probability of success under an ordered decision process}
\label{sec:p15-ssat-purpose}

Stochastic Boolean satisfiability is not SAT with randomly chosen
heuristics, and not merely a probability attached to a propositional
formula. Its prefix specifies which Boolean choices are controlled,
which are randomized, and what information is available before each
choice. The task is to calculate the greatest probability of satisfying
the matrix using a legal contingent policy, or to decide whether that
value meets a specified threshold.

This presentation uses the directly read Littman--Majercik--Pitassi
January 2000 author manuscript, including its algorithms and appendices
\cite{LittmanMajercikPitassi2000SSATAuthor}. The final journal article,
Majercik's 2009 handbook chapter, and the 2021
handbook chapter are distinct source presentations. The 2002 erratum corrects author
affiliations, not the mathematics
\cite{LittmanMajercikPitassi2002SSATErratum}. The definitions below distinguish
the native fair-coin presentation from the declared independent biased extension.

\begin{hbdefinition}{def:p15-ssat-native}{The native ordered map and random-quantifier glyph}
The original input is \((\phi,\theta,Q)\), with ordered Boolean
variables \(x_1,\ldots,x_n\) and
\(Q(x_i)\in\{\exists,\forall,\ssatrandom\}\).
The native symbol \(\ssatrandom\) is a Roman R rotated by 180 degrees;
it is neither an ordinary upright R nor an ordinary first-order quantifier.
Its source interpretation is a fair independent Boolean draw.
The original restriction is \(\phi\upharpoonright_{x=b}\), including
order-preserving renumbering of remaining variables and their quantifier map.
Its value is written \(\operatorname{val}(\phi,Q)\). At a nonterminal
whose outermost variable is \(x_1\), it combines the two restricted
values by maximum for \(\exists\), minimum for \(\forall\), and
their arithmetic mean for \(\ssatrandom\); an empty clause has value
zero and an empty matrix has value one. These are the five native clauses;
the positive native instance satisfies \(\operatorname{val}(\phi,Q)\geq\theta\)
\cite[Table~I and pp.~3--4]{LittmanMajercikPitassi2000SSATAuthor}.
\end{hbdefinition}

\begin{hbdefinition}{def:p15-ssat-input}{Formula, ordered quantifiers and threshold}
A closed finite stochastic input is \((\phi,\theta,\Pi)\), where
\(\phi\) is a Boolean CNF, every matrix variable occurs once in the
finite scalar prefix \(\Pi=Q_1x_1\cdots Q_nx_n\), and
\(\theta\in[0,1]\). The \(x_i\) are distinct; unused prefix
variables and an empty prefix are allowed. Quantifiers are existential
\(\exists\) or randomized \(\ssatrandom^{p_i}\); the source's
\emph{Extended SSAT} additionally allows universal \(\forall\).
The superscript \(p\) annotates the native random glyph in the declared
biased extension; dropping it denotes the source's fair case.

The original's main random quantifier is fair, \(p=1/2\). In the
declared biased extension, each \(p_i\in[0,1]\) is fixed in advance,
and the random variables are independent Bernoulli draws with
\(P(x_i=1)=p_i\). Exact computational inputs encode these probabilities
and \(\theta\) as binary rationals. A mathematical recurrence may
also use specified real probabilities, without thereby claiming a finite
exact real-arithmetic implementation.

Literals are \(x\) and \(\bar x\); the bar is Boolean negation.
Clauses are finite sets of literals. An empty clause is false and an
empty set of clauses is true. Restriction \(\phi|_{x=b}\) deletes
clauses containing a true literal, removes false literals from remaining
clauses, and retains any resulting empty clause. The corresponding
quantifier is removed and the relative order of all others retained.
There are no free variables, identity predicates or unspecified random
dependencies in this input language.
\end{hbdefinition}

The source writes its quantifier map as \(Q(x_i)\), with indices
ordered from outermost to innermost. We retain that order through the
prefix rather than renumbering away the dependency information
\cite[sec.~1.2.1, pp.~3--5]{LittmanMajercikPitassi2000SSATAuthor}.

The prefix representation below is comparison data, not a replacement for
the native map notation: \(\Pi_Q=Q(x_1)x_1\cdots Q(x_n)x_n\).
In the fair case a native random occurrence receives parameter \(1/2\).
The restriction notation \(\phi|_{x=b}\) used in the algorithms denotes
the simplified Boolean cofactor with old names retained; the native
renumbering is an order-preserving renaming of that object.

\begin{hbdefinition}{def:p15-ssat-value}{Value and the decision question}
The value \(V(\Pi,\phi)\) is recursively determined. Empty-clause
and empty-matrix terminals have values zero and one. Otherwise, for
the outermost variable \(x\), let \(v_b\) be the value after its
restriction to \(b\). Existential quantification combines the children
by \(\max(v_0,v_1)\); universal quantification combines them by
\(\min(v_0,v_1)\); random quantification combines them by
\((1-p)v_0+pv_1\). A positive decision instance satisfies
\(V(\Pi,\phi)\geq\theta\). At \(\theta=0\), every well-formed
input is positive, even a contradictory matrix. This is not the SAT
question of whether there exists one satisfying valuation.
\end{hbdefinition}

\section{The evaluation calculus and what its strategies may read}
\label{sec:p15-ssat-evaluation}

\begin{hblemma}{thm:p15-ssat-native-correspondence}{proved=annex:p15.ssat-native-correspondence}{Exact fair-source and prefix-value correspondence}
For the native fair input and its order-preserving prefix encoding,
\(\operatorname{val}(\phi,Q)=V(\Pi_Q,\phi)\).
Order-preserving variable renaming together with its transported
quantifier map preserves either value, and hence the decision at the same
threshold. This does not cross quantifier positions or identify a biased
draw with a fair one.
\end{hblemma}

\begin{hbdefinition}{def:p15-ssat-evaluation-judgment}{The numerical evaluation judgment}
For a well-formed prefix and matrix, \((\Pi,\phi)\Downarrow v\)
is the value-evaluation judgment. Its semantic interpretation is
\(V(\Pi,\phi)=v\); a derivation uses the five rules on the following
sheet. Thus \(v\) is a real continuation value, not a Boolean valuation
of every variable. The sheet computes the fair source case and the explicitly
specified independent weighted extension under their respective parameters.
\end{hbdefinition}

\calculussheet{sheet:p15-stochastic-satisfiability}{Finite SSAT value evaluation}{Language: closed finite prenex Boolean CNF; judgment $(\Pi,\phi)\Downarrow v$; literal complement $\bar x$ is Boolean negation; no equality predicate. Fair source case $p=1/2$, declared independent biased extension $0\leq p\leq1$.}
\hbrule{rule:p15-ssat-false}{empty-clause}{}{(\Pi,\phi)\Downarrow0}{$\phi$ contains an empty clause}
\hbrule{rule:p15-ssat-true}{empty-matrix}{}{(\Pi,\varnothing)\Downarrow1}{\none}
\hbrule{rule:p15-ssat-exists}{max}{(\Pi',\phi|_{x=0})\Downarrow v_0\quad(\Pi',\phi|_{x=1})\Downarrow v_1}{(\exists x\Pi',\phi)\Downarrow\max(v_0,v_1)}{$x$ is the outermost prefix variable; restrictions preserve the rest of the prefix}
\hbrule{rule:p15-ssat-forall}{min}{(\Pi',\phi|_{x=0})\Downarrow v_0\quad(\Pi',\phi|_{x=1})\Downarrow v_1}{(\forall x\Pi',\phi)\Downarrow\min(v_0,v_1)}{Extended SSAT; $x$ is outermost}
\hbrule{rule:p15-ssat-random}{weighted}{(\Pi',\phi|_{x=0})\Downarrow v_0\quad(\Pi',\phi|_{x=1})\Downarrow v_1}{(\ssatrandom^p x\Pi',\phi)\Downarrow(1-p)v_0+pv_1}{$x$ is outermost; $p$ is its fixed independent Bernoulli parameter}


\begin{hbdefinition}{def:p15-ssat-strategy}{Legal contingent strategies}
A deterministic existential strategy assigns each controlled variable
from the already revealed prefix history, never from a later random
or universal value. In an existential/random prefix this can be written
as one Boolean table for each existential variable, whose arguments are
the preceding random values; its own earlier choices are recovered from
its earlier tables. In Extended SSAT the opposing universal player may
likewise read preceding revealed history but not later choices or draws.
The complete strategies specify values even on zero-probability histories.

For a fixed pair of strategies, a full assignment to the random variables
induces one complete play. Its weight is the product of their specified
Bernoulli outcome weights. Success is the indicator that the final matrix
is true. The expected success is the finite weighted sum of those
indicators. Without universal variables, the objective maximizes that
expectation. With universal variables it is a finite perfect-information
zero-sum game against a minimizing player and independent chance.
\end{hbdefinition}

\begin{hbtheorem}{thm:p15-ssat-strategy-semantics}{proved=annex:p15.ssat-strategy-semantics}{Complete evaluation and legal optimal strategies}
The displayed finite calculus is sound and complete for the recursive
value. Its value lies in \([0,1]\), terminates after finitely many
prefix restrictions, and equals the optimal expected success under the
declared strategy information. In the extended game, deterministic
backward strategies achieve matching maximizing/minimizing guarantees.
Randomizing a controlled choice cannot improve these guarantees.
The result includes empty prefixes, unused variables, deterministic
random probabilities zero/one, and histories of zero probability.
\end{hbtheorem}

Values at a zero-probability history are defined continuation-game values,
not conditional probabilities obtained by dividing by that history's
probability. A zero-weight child contributes nothing to its random parent,
although its continuation game remains syntactically well defined.
Strategy tables may differ arbitrarily there without changing the root
value. This distinction uses the null-event boundary already established
in Chapter~\ref{sec:p15-finite-probability}.

\begin{hbtheorem}{thm:p15-ssat-prefix-order}{proved=annex:p15.ssat-prefix-order}{Commuting like operators and crossing an information boundary}
Adjacent existential quantifiers commute, adjacent universal quantifiers
commute, and adjacent independent random quantifiers commute with their
own fixed biases retained. In contrast,
\[
 V(\exists x\,\ssatrandom^p y,\ x\leftrightarrow y)=\max(p,1-p),
 \qquad
 V(\ssatrandom^p y\,\exists x,\ x\leftrightarrow y)=1.
\]
For every two-variable continuation payoff \(f(a,b)\in[0,1]\),
\[
 \max_a\sum_b\mu_bf(a,b)
       \leq\sum_b\mu_b\max_af(a,b).
\]
Thus moving a choice after an independent random observation relaxes
the information restriction and can increase value. It is not generally
a value-preserving variable-order heuristic.
\end{hbtheorem}

At \(p=1/4\), the crossed-prefix values are \(3/4\) and one;
threshold \(4/5\) distinguishes them. The policy \(x=y\) is legal
only in the second prefix. In the first, the controlled value must be
constant. This is a counterexample about order, not uncertainty in the
matrix's propositional truth table.

\begin{hbexample}{def:p15-ssat-source-three-variable}{The source's ordered three-variable calculation}
The original's matrix is
\[
 \phi=(x\lor y)\land(y\lor\bar z)\land(\bar x\lor\bar y).
\]
For fair prefix \(\ssatrandom^{1/2}x\,\exists y\,\ssatrandom^{1/2}z\),
the \(x=0\) branch chooses \(y=1\) and has value one;
the \(x=1\) branch must choose \(y=0\) and succeeds precisely
when \(z=0\), so has value \(1/2\). The root is \(3/4\).
Moving \(z\) before \(y\) retains \(3/4\): only the pair
\(x=z=1\) cannot be satisfied by any \(y\). Moving \(y\)
before both draws instead gives value \(1/2\): constant \(y=1\)
succeeds for \(x=0\), while constant \(y=0\) succeeds only for
\(x=1,z=0\). Thus splitting on the later choice as though it were
an outermost variable computes a different problem
\cite[pp.~4,37]{LittmanMajercikPitassi2000SSATAuthor}.
\end{hbexample}

\section{Source DPLL simplifications and their exact scope}
\label{sec:p15-ssat-dpll}

\begin{hbtheorem}{thm:p15-ssat-unit-pure}{proved=annex:p15.ssat-unit-pure}{Unit and pure-literal transformations at arbitrary prefix positions}
Let \(\Pi\setminus x\) remove a variable while retaining every
other quantifier and its order. If the matrix contains the unit clause
forcing \(x=b\), then
\[
 V(\Pi,\phi)=
 \begin{cases}
 V(\Pi\setminus x,\phi|_{x=b}),&Q(x)=\exists,\\
 0,&Q(x)=\forall,\\
 p_bV(\Pi\setminus x,\phi|_{x=b}),&Q(x)=\ssatrandom^p,
 \end{cases}
 \qquad p_1=p,\quad p_0=1-p.
\]
This holds even when \(x\) is not outermost, because the opposite
branch has success value zero and all preceding operators are positively
homogeneous. If \(p_b=0\), the value is zero without evaluating a
division by \(p_b\).

If a variable occurs only with sign \(b\), an existential variable
may be fixed to \(b\) and a universal variable to \(1-b\), at
any prefix position. A pure random variable cannot generally be fixed
to either sign without changing the value. Tautological clauses can
be removed and unused variables omitted without changing value.
\end{hbtheorem}

The original calls these operations unit resolution and purification.
They are exceptional simplifications justified by matrix structure, not
permission to cross an arbitrary quantifier boundary. For example the
pure random one-clause input \(\ssatrandom^{1/4}x\;(x)\) has
value \(1/4\), not one. A later random unit with a history-dependent
probability would not have the constant factor used in the theorem;
that is outside the independent-input hypothesis
\cite[sec.~2.1, pp.~14--16]{LittmanMajercikPitassi2000SSATAuthor}.

\begin{hbdefinition}{def:p15-ssat-window-contract}{The exactness contract for a threshold window}
For real thresholds \(L\leq H\), an evaluator's return \(u\)
represents true value \(v\) by the following contract:
\[
 L\leq v\leq H\ \Longrightarrow\ u=v,\qquad
 v<L\ \Longrightarrow\ u<L,\qquad
 v>H\ \Longrightarrow\ u>H.
\]
The contract says that the returned number outside the window need not
be the exact value; it must remain on the correct side. It is not a
probability interval around an approximate numerical estimate.
\end{hbdefinition}

\begin{hbmethod}{def:p15-ssat-window-evaluator}{A corrected source-style thresholded evaluator}
Define \(\operatorname{Eval}(\Pi,\phi,L,H)\) as follows, using
exact arithmetic and restrictions. Return the exact terminal zero or
one when an empty clause or empty matrix is detected. Apply any unit
or pure-literal transformation from the preceding theorem. A random
unit of weight \(s>0\) returns
\(s\operatorname{Eval}(\Pi\setminus x,\phi|_{x=b},L/s,H/s)\);
weight zero returns zero. Other such transformations leave thresholds
unchanged.

If none applies, split on an outermost variable. Write
\(\operatorname{Eval}_b(L',H')\) for the recursive call on its
cofactor and remaining prefix.
\begin{enumerate}
\item Existential: obtain \(u_0=\operatorname{Eval}_0(L,H)\).
      If \(u_0>H\), return \(u_0\). Otherwise obtain
      \(u_1=\operatorname{Eval}_1(\max(L,u_0),H)\), and return
      \(\max(u_0,u_1)\).
\item Universal: obtain \(u_0=\operatorname{Eval}_0(L,H)\).
      If \(u_0<L\), return \(u_0\). Otherwise obtain
      \(u_1=\operatorname{Eval}_1(L,\min(H,u_0))\), and return
      \(\min(u_0,u_1)\).
\item Random with \(p=0\) or \(p=1\): evaluate only the
      positive-weight cofactor with the unchanged window.
\item Random with \(0<p<1\): let \(q=1-p\) and obtain
      \(u_0=\operatorname{Eval}_0((L-p)/q,H/q)\).
      If \(qu_0+p<L\), return \(qu_0\). If \(qu_0>H\),
      return \(qu_0\). Otherwise obtain
      \(u_1=\operatorname{Eval}_1((L-qu_0)/p,(H-qu_0)/p)\)
      and return \(qu_0+pu_1\).
\end{enumerate}
Recursive windows may extend beyond \([0,1]\); that is legitimate.
Every recursive call removes a variable. Ties are handled by the
strict early-return inequalities, not by silently changing the decision
threshold from \(\geq\) to \(>\).
\end{hbmethod}

At \(p=1/2\), the transformed random windows are
\((2L-1,2H)\) and \((2L-u_0,2H-u_0)\), precisely the
source's scaling structure. The explicit corrections are strict high
early-return tests and zero/one-biased endpoint handling. The proof
does not certify a literal transcription of the original Figure 2.

\begin{hbtheorem}{thm:p15-ssat-window-correctness}{proved=annex:p15.ssat-window-correctness}{Every window, every quantifier and every probability endpoint}
The corrected evaluator terminates and satisfies the declared window
contract for every finite valid input and real \(L\leq H\).
Calling it with \((L,H)=(0,1)\) gives the exact value; calling it
with \((\theta,\theta)\) decides the inclusive threshold by comparing
its return with \(\theta\). The result includes arbitrary-position
unit/pure transformations and all biased random endpoint cases.
\end{hbtheorem}

\begin{hbexample}{def:p15-ssat-source-endpoint-failure}{Why the printed closed-window claim needs correction}
The source returns immediately when an existential first child meets
\(u_0\geq H\), while Appendix A claims exactness throughout the
closed window \([L,H]\). In the three-variable source example
\(\exists y\,\ssatrandom^{1/2}x\,\ssatrandom^{1/2}z\;\phi\)
above, the \(y=0\) child has value \(1/4\), the \(y=1\) child
has value \(1/2\). With \(L=0,H=1/4\), returning \(1/4\)
immediately need not report the true value \(1/2\); it also fails
the requirement to return strictly above \(H\) for values above it.
The matrix has no unit or pure variable, so these simplifications do
not eliminate the countercase. The corrected strict test evaluates the
second child. The original Appendix A random argument also prints a
division by four where the defining mean divides by two; the proof here
uses the actual weighted recurrence
\cite[Fig.~2, p.~15; Appendix~A, pp.~48--49]{LittmanMajercikPitassi2000SSATAuthor}.
\end{hbexample}

\section{Exact certificates and computational guarantees}
\label{sec:p15-ssat-certificates}

\begin{hbmethod}{def:p15-ssat-value-certificate}{A full finite value certificate}
A certificate is a canonical prefix-respecting assignment tree.
Each node records its remaining prefix, restricted matrix and rational
value. A terminal supplies the actual empty-clause or empty-matrix
condition; a split supplies both child restrictions and the indicated
max, min or weighted combination. Check the original prefix, all
restrictions, terminal conditions and rational equations, then compare
the root value to the threshold. A zero-weight child may still be
supplied; no undefined null-event conditional is checked. Arbitrary
matrix-SAT models or trees that cross prefix boundaries are not these
certificates.
\end{hbmethod}

\begin{hbtheorem}{thm:p15-ssat-certificate-complexity}{proved=annex:p15.ssat-certificate-complexity}{Certificate correctness and polynomial-space decision}
The stated checker terminates and accepts precisely correct canonical
value trees. Every valid input has such a finite tree, so exact threshold
decisions have complete certificates. Checking is polynomial in the
explicit input and certificate length with binary rational arithmetic;
the certificate itself may be exponential.

For binary rational independent weights, exact depth-first evaluation
uses polynomial space. The existential/random-only threshold problem
is \(\mathsf{PSPACE}\)-complete: replacing each universal QBF
quantifier by a fair random quantifier and using threshold one is a
polynomial reduction. This completeness is for unbounded alternation
and inclusive threshold; it does not classify every fixed-prefix subtype.
\end{hbtheorem}

The QBF prerequisite is the actual finite Boolean language and TQBF
completeness result in Chapter~\ref{sec:p04-complexity-complete-problems};
the prefix-relative certificate discipline is also developed in
Chapter~\ref{sec:p12-q-resolution-and-quantified-certificates}.
An SSAT value certificate is not a Q-resolution refutation: it certifies
a numerical optimization/expectation, rather than a false QBF.

The source separately discusses DPLL-tree search. Its allowed search
order preserves blocks, with justified unit/pure exceptions, and its
universal search dovetails calls for literals in the current block.
When the fastest cofactor finishes, it retains that cofactor tree and
constructs the opposite cofactor tree, then joins them at the selected
variable. Its reported \(m n^{O(\log s)}\) bound uses the smallest
\emph{allowed DPLL tree} size \(s\), not an arbitrary matrix decision
tree that ignores the prefix. That proof/size-analysis and the cited
decision-tree lower bounds remain separately attributed source results;
neither is implied by the finite evaluator's termination theorem
\cite[sec.~2.2, pp.~18--20]{LittmanMajercikPitassi2000SSATAuthor}.

\section{Weighted model counting is a special prefix, not the entire calculus}
\label{sec:p15-ssat-counting}

\begin{hbtheorem}{thm:p15-ssat-projections}{proved=annex:p15.ssat-projections}{The exact random-only, exist-random and random-exist cases}
Let \(R\) be independent random variables with product mass \(\mu\)
and let \(X\) be controlled variables.
\begin{enumerate}
\item A random-only prefix has value
      \(\sum_r\mu(r)\mathbf1_{\phi(r)}\). For \(n\) fair
      variables this is \(\#\mathrm{models}(\phi)/2^n\).
\item An exist-random prefix has value
      \(\max_x\sum_r\mu(r)\mathbf1_{\phi(x,r)}\), with one
      fixed \(x\) selected before any observation.
\item A random-exist prefix has value
      \(\sum_r\mu(r)\mathbf1_{\exists x\phi(x,r)}\), allowing
      a separately chosen satisfying extension for each observed \(r\).
\end{enumerate}
The last case is a weighted projected count, not a count of every
\((x,r)\) satisfying extension. Existential auxiliary variables placed
after all random variables preserve this value exactly when their
projection defines the same event on the random assignments.
Randomly quantifying the same auxiliaries can introduce multiplicity
weights and change it.
\end{hbtheorem}

For a fresh unconstrained Boolean auxiliary, adding it existentially
after a random-only formula preserves its value. Randomly adding it
also preserves value when truly unconstrained, since its weights sum to
one; but forcing that auxiliary to equal a deterministic gate output
and then assigning it a fair random quantifier introduces a factor
\(1/2\). With \(g\) such uniquely constrained gate outputs, the
factor is \(2^{-g}\). SAT equisatisfiability alone therefore is not a
proof of probability preservation.

\section{An explicit Bayesian-network inference translation}
\label{sec:p15-ssat-bayesian-inference}

\begin{hbdefinition}{def:p15-ssat-row-coin-encoding}{Boolean DAG, row coins and deterministic constraints}
Take a finite acyclic Boolean Bayesian network in topological order
\(A_1,\ldots,A_m\), with explicit rational conditional tables
\(t_{i,u}=P(A_i=1\mid\operatorname{pa}(A_i)=u)\).
Each table row has a total probability convention, including rows whose
parent event may be null. Introduce one independent Bernoulli coin
\(C_{i,u}\) of probability \(t_{i,u}\) per row. Put all those
coins first in a random prefix. Then existentially quantify all \(A_i\)
and any Boolean gate auxiliaries, constrained by full gate equivalences
and
\[
 A_i\leftrightarrow
       \bigvee_u[\operatorname{Pattern}_{i,u}(\operatorname{pa}(A_i))
                         \land C_{i,u}].
\]
The parent patterns are disjoint and exhaustive, including the unique
empty pattern for a root node. Let \(\Psi\) be the conjunction of
these constraints. Attach any Boolean evidence/query event in the
\(A_i\)'s to that matrix. Gate equivalences are converted to CNF
with their auxiliaries in the final existential block.
Write \(\boldsymbol C\) for the row coins in the chosen prefix order,
and \(\boldsymbol t\) for their probabilities in that same order.
Thus \(\ssatrandom^{\boldsymbol t}\boldsymbol C\) is the corresponding
weighted block, not a fair block inferred from omission of individual indices.
\end{hbdefinition}

\begin{hbtheorem}{thm:p15-ssat-bn-translation}{proved=annex:p15.ssat-bn-translation}{Exact joint probability and qualified conditional inference}
Every coin assignment has exactly one \(A\)-assignment satisfying
\(\Psi\), and every gate auxiliary is uniquely determined. The
induced law of \(A\) is exactly the Bayesian-network joint law.
Therefore, for any event \(B(A)\),
\[
 V(\ssatrandom^{\boldsymbol t}\boldsymbol C\,\exists\boldsymbol A,\Psi\land B)
      =P_{\mathrm{BN}}(B).
\]
In particular evidence \(E\) and query \(H\) give
\[
 P_{\mathrm{BN}}(H\mid E)
    =\frac{V(\ssatrandom^{\boldsymbol t}\boldsymbol C\,\exists\boldsymbol A,
                    \Psi\land H\land E)}
            {V(\ssatrandom^{\boldsymbol t}\boldsymbol C\,\exists\boldsymbol A,
                    \Psi\land E)}
\]
only when the denominator is positive. At zero evidence probability,
the ratio is undefined. The encoding is polynomial in the explicit
table and event-circuit size, not necessarily in a succinct description
of exponentially many CPT rows.
\end{hbtheorem}

The constrained existential variables are logical auxiliaries, not free
actions. Their full equivalences leave no freedom to change the network
realization. For planning, an ordinary SSAT choice may read every earlier
random bit: placing a hidden state draw before an action can leak
information unless an encoding proves observation-only correspondence
or supplies a separately justified dependency restriction. Conditioning
on evidence also is not intervention on a network mechanism.

The earlier original connects belief-network inference with randomized
satisfiability on pp.~5--8. This explicit row-coin construction states
and proves the encoding conditions rather than treating that connection
as a general equivalence for every mixed prefix. It is an independent
construction; it is not attributed as a literal source encoding with
these auxiliary variables.

\begin{hbexample}{def:p15-ssat-bn-worked}{Evidence requires a denominator and dependence requires an encoding}
Let \(P(X=1)=1/5\), \(P(Y=1\mid X=1)=4/5\), and
\(P(Y=1\mid X=0)=1/10\). Use three independent row coins and
constraints \(X\leftrightarrow C_X\),
\(Y\leftrightarrow[(X\land C_{Y,1})\lor(\bar X\land C_{Y,0})]\).
The evidence matrix with \(Y\) has value
\[
 \frac15\frac45+\frac45\frac1{10}=\frac6{25},
\]
and the query/evidence matrix with \(X\land Y\) has value
\(4/25\). Their ratio is \(2/3\). Using the marginal probabilities
of \(X,Y\) as independent random quantifiers would instead give
\(P(X\land Y)=(1/5)(6/25)=6/125\), hence conditional \(1/5\):
it destroys the dependency, despite using both correct marginals.
\end{hbexample}

\section{Policy trees, empirical optimization and incomplete local search}
\label{sec:p15-ssat-policy-sampling}

The original's sampling method concerns existential/random prefixes,
not the adversarial universal extension. Its policy tree gives one copy
of an existential variable for each preceding random history. A policy
assigns those copies, and its value is the weighted proportion of
random leaves whose path formula is satisfied. Sampled histories must
share an earlier controlled decision whenever their earlier random
histories agree; giving a separate early decision to every full sample
would introduce clairvoyance.

\begin{hbdefinition}{def:p15-ssat-sample-policy}{A complete policy and a sampled policy-tree value}
If existential variable \(x_i\) has \(r_i\) preceding random
variables, its full table has \(2^{r_i}\) entries. Put
\(D=\sum_{i:Q_i=\exists}2^{r_i}\); there are \(P=2^D\)
labelled deterministic table policies, including values on null histories.
Draw \(w\geq1\) independent full random assignments from the declared
product mass, independently of a policy under evaluation. For a policy
\(s\), its empirical value \(\widehat v_s\) is the fraction of
the \(w\) induced plays satisfying the matrix. The empirical optimum
is \(\widehat v_* = \max_s\widehat v_s\).

The source's \texttt{sampleevalssat} splits the current sample multiset
according to an outer random variable and weights the children by
their multiplicities. An empty child contributes zero and is not
evaluated through a ratio \(0/0\). At an outer existential node,
maximize the two restricted subproblems using the \emph{same} sample
histories. Recurse to the empty-matrix/empty-clause terminals. This
computes the exact optimum on the sampled tree, not the exact optimum
under the original distribution. Any unvisited policy-table entries
must be completed arbitrarily to obtain a total legal policy.
\end{hbdefinition}

\begin{hbtheorem}{thm:p15-ssat-sampling-bound}{proved=annex:p15.ssat-sampling-bound}{Sampling correctness and the cost of choosing a policy on the sample}
For any fixed complete policy, \(\widehat v_s\) is unbiased and
\[
 P(|\widehat v_s-v_s|>\varepsilon)
       \leq2e^{-2w\varepsilon^2},\qquad\varepsilon>0.
\]
Uniformly over the finite policy collection,
\[
 P(\exists s:|\widehat v_s-v_s|>\varepsilon)
       \leq2P e^{-2w\varepsilon^2}.
\]
On the complementary event,
\(|\widehat v_*-V|\leq\varepsilon\), and a sample-optimal complete
policy has true value at least \(V-2\varepsilon\).
Thus \(w\geq\log(2P/\delta)/(2\varepsilon^2)\) suffices for
these simultaneous guarantees with probability at least \(1-\delta\).
The recursive sampled-tree procedure terminates and computes
\(\widehat v_*\) exactly with the stated empty-child convention.
\end{hbtheorem}

This is the source's fixed-policy versus all-policies distinction
\cite[sec.~2.3, Lemma~4 and Theorem~5, pp.~21--30]{LittmanMajercikPitassi2000SSATAuthor}.
For random-only inputs \(P=1\); for an outer block of \(c\) choices,
\(P=2^c\). With alternating choices, \(D\), and therefore
\(\log P\), can be exponential. A fixed-policy bound independent
of input size is not a polynomial guarantee for optimizing all policies.

\begin{hbtheorem}{thm:p15-ssat-sampling-overfit}{proved=annex:p15.ssat-sampling-overfit}{The source's exponential sampling failure case}
Use fair prefix
\[
 \ssatrandom^{1/2}y_1\cdots\ssatrandom^{1/2}y_r\;
 \exists x\;\ssatrandom^{1/2}z\;(x\leftrightarrow z).
\]
Its true value is \(1/2\), but it has \(2^{2^r}\) contingent
policies. If \(N_{h,b}\) counts samples with initial history
\(h\in\{0,1\}^r\) and final bit \(b\), its empirical optimum
is
\[
 \widehat v_* =\frac1w\sum_h\max(N_{h,0},N_{h,1}).
\]
It satisfies
\(E[\widehat v_*]\geq1-(w-1)/2^{r+1}\).
For \(0<\varepsilon<1/2\), a guarantee
\(P(|\widehat v_*-1/2|\leq\varepsilon)\geq1-\delta\)
therefore requires
\[
 w\geq1+2^r(1-2\varepsilon)(1-\delta).
\]
For example \(r=10,\varepsilon=1/10,\delta=1/20\) requires
at least 780 samples. This is a necessary bound, not a sufficient one.
\end{hbtheorem}

The original Appendix B uses \(r=n-2\). The normalized majority-count
formula here states exactly what the empirical optimizer computes;
the source's match-sum display is not copied as a probability without
its normalization and majority selection. The failure is optimization
over many scarcely observed histories, not inaccurate evaluation of
one policy fixed independently of the sample
\cite[Appendix~B, pp.~49--50]{LittmanMajercikPitassi2000SSATAuthor}.

The source's \texttt{randevalssat} additionally replaces exact sampled
optimization by hillclimbing. Each sampled path yields a Boolean path
formula over the legally shared decision copies. Its two-level objective
counts satisfied clauses and satisfied entire path formulas, weighting
a whole path by the number of clauses of the original matrix. It starts
from random assignments, chooses an unsatisfied clause, flips a variable
when the objective improves, otherwise considers a least-damaging flip
with probability one half, and restarts after stipulated failure,
iteration and restart limits. More satisfied scattered clauses need not
mean more satisfied complete paths.

If a returned policy is empirically \(g\) below the sampled optimum,
the uniform sampling event only guarantees true value at least
\(V-g-2\varepsilon\). Without a proved \(g\), the sampling theorem
does not certify this local search as globally optimal. The source
explicitly reports local optima, missing policy branches and memory
cost as limitations. It does not supply a finite-error decision proof
for every returned approximate number or an extension of this sampling
routine to all universal/minimizing prefixes.

\section{The random-instance distribution is not random quantification}
\label{sec:p15-ssat-random-instances}

The original's fixed-clause model \(F^ {k,n}_m\) first draws a CNF
instance: each of \(m\) independently generated clauses chooses \(k\)
distinct variables uniformly, then independently chooses each sign with
probability one half. This outer distribution over formulas is different
from the inner random quantifiers of one SSAT instance. Any claimed
average solver improvement must identify which distribution and which
prefix/threshold settings were used.

\begin{hbtheorem}{thm:p15-ssat-random-formula-bound}{proved=annex:p15.ssat-random-formula-bound}{The first-moment bound with its actual asymptotic margin}
For \(n\geq k\geq1\), \(m\geq0\), and integer \(0\leq t\leq n\),
let \(Z\) count satisfying assignments of a fixed-clause random formula.
Then
\[
 E[Z]=2^n(1-2^{-k})^m,\qquad
 P(Z\geq2^t)\leq2^{n-t}(1-2^{-k})^m.
\]
For \(k=3\), the latter tends to zero whenever
\[
 m\log(8/7)-(n-t)\log2\longrightarrow+\infty.
\]
Equality at \(m=(n-t)\log2/\log(8/7)\) alone gives bound one,
not an almost-sure negative classification. Under the stated distinct-
variable clause model, the expected number of unused variables is
\(n(1-k/n)^m\).
\end{hbtheorem}

The original discusses first-moment thresholds on pp.~9--12. Its
unused-variable expression \(n(1-1/n)^{km}\) instead corresponds
to independent variable choices with replacement; it is not the exact
value for clauses whose variables are distinct. Both approximate
\(ne^{-km/n}\) in appropriate sparse regimes, but the exact sampling
convention and the necessary asymptotic margin remain visible here.
No sharp-threshold theorem or empirical hardness claim follows from
this first-moment calculation alone.

\section{What the original experiments do and do not establish}
\label{sec:p15-ssat-experiments}

The primary paper tests its algorithms on small random instances, not
on a distribution of all possible probabilistic inference problems.
Its first comparisons use 1000 formulas with 20 variables, 3 literals per
clause and 141 generated clauses, taking prefixes of that clause list.
Thresholds use \(\theta=2^{t-n}\). Exist-random experiments vary
the number of initial choices, clause count and threshold; the observed
peak work does not follow the ordering of worst-case complexity classes.
These work measurements count recursive calls, not a theorem that one
prefix class is universally faster than another.

The heuristic experiment uses eight sets of 100 formulas: 24 variables
with 24, 48, 72, 96 clauses and 32 variables with 32, 64, 96, 128 clauses, all
3-CNF. Its modified generator rejects formulas until all variables
occur, changing the unconditioned fixed-clause distribution. Prefixes
alternate equal-size blocks and begin either existentially or randomly.
The source compares RAND, SATS MOST, 2S JW, POS 2S JW, MOMS and MAX UNIT,
always selecting within the current outer block. The 10,400 instances
were solved by a C implementation on a Sun Ultra-1, SunOS-5.7,
128 MB RAM; the metric is change in recursive-call count relative to
strict prefix-order splitting. The author qualifies transfer to larger
or structured planning instances
\cite[sec.~3.3, pp.~36--41]{LittmanMajercikPitassi2000SSATAuthor}.

The sampling/local-search experiment uses 27 sets of 100 formulas from
the cross product \(n\in\{8,12,16\}\),
\(m\in\{n,2n,3n\}\), \(k\in\{3,4,5\}\), again conditioned
on all variables appearing. It constructs 25,200 prefix instances and
ten sampled trees per instance, with sample sizes ranging from 5 to 4086
on an approximately logarithmic scale. Performance is measured by
squared error against exact full-tree values; local-search iteration
and restart budgets grow with the constructed problem. More samples
remove sampling error for exact empirical optimization, but do not
remove local-search error. The reported anomalous 12-variable, 24-clause,
3-CNF experiment with block size 4 and an initial existential block
has unequal random/choice counts; its squared-error behavior is not
pooled silently with the balanced-prefix experiments
\cite[sec.~3.4, pp.~42--45]{LittmanMajercikPitassi2000SSATAuthor}.

\begin{systementry}{entry:p15-stochastic-satisfiability}{Finite stochastic Boolean satisfiability}
\entrylanguage{Closed finite ordered Boolean CNF with existential and fixed independent rational random quantifiers; universal min belongs to the explicitly declared Extended SSAT. Literal complement is Boolean negation; no identity predicate.}
\entrysemantics{Optimal expected matrix satisfaction under nonanticipating prefix strategies; maximizing/minimizing perfect-information continuation games in the extended case.}
\entrycalculus{sheet:p15-stochastic-satisfiability}
\entryadequacy{proved=annex:p15.ssat-strategy-semantics}{The finite value rules exactly compute the declared sequential-strategy value, including null histories.}
\entrydecidability{proved=annex:p15.ssat-certificate-complexity}{Exact rational threshold decision terminates, has complete value certificates and is PSPACE-complete for unbounded existential/random alternation.}
\entrycharacteristic{Crossing random/choice boundaries changes information and can change value; random-only counting, contingent projected counting and conditional Bayesian inference require their distinct stated conditions. Sampling optimization and bounded local search are not exact decision certificates.}
\entrysources{LittmanMajercikPitassi2000SSATAuthor,LittmanMajercikPitassi2002SSATErratum}
\end{systementry}

\begin{proofboundary}{def:p15-ssat-source-boundary}{Exact value, information and algorithmic scope}
The annex proves the declared value/strategy calculus, prefix boundaries,
source-style simplifications and corrected threshold evaluator,
certificate/space guarantees, special-prefix projections, row-coin
inference construction and sampling/first-moment bounds. It does not
identify the source's fair presentation with every weighted or dependency-
restricted stochastic constraint language.
The source's universal-search optimal-tree size analysis, full cited
random-formula lower bounds and subtype approximation-complexity
results retain their separately cited standing above. The 2001 final journal body,
2009 handbook chapter and 2021 chapter remain distinct source presentations;
stochastic constraint programming and all later chapter-specific
constructions and experiments require their own defined inputs, information
constraints, algorithms and correctness statements.
\end{proofboundary}
